% 12
\begin{frame}[t]{本课题的实验条件：起步阶段与合作阶段}
\vspace{.05cm}
\centering
\begin{tikzpicture}[st/.style={box,inner sep=5pt,text width=2.05cm,minimum height=1.0cm,font=\scriptsize},
  lt/.style={font=\small\bfseries,text=DarkBlue,anchor=west}]
 \node[lt] (l1) at (0,0) {\faSeedling\enspace 起步阶段 · 公开资源，不依赖医院数据};
 \node[st,anchor=west] (a1) at (0,-.95) {\ico{\faDatabase}\ \textbf{公开数据}\\Perturb-seq、VCC 2025};
 \node[st,right=.3cm of a1] (a2) {\ico{\faFilter}\ \textbf{审计与预处理}\\规模、ID 映射、划分};
 \node[st,right=.3cm of a2] (a3) {\ico{\faRuler}\ \textbf{基线}\\CPU 即可完成};
 \node[st,right=.3cm of a3,fill=PaleTeal,draw=DarkBlue,line width=.9pt] (a4) {\ico{\faLock}\ \textbf{评分脚本}\\实验前冻结};
 \node[st,right=.3cm of a4] (a5) {\ico{\faMicrochip}\ \textbf{公开预测器}\\单卡 24--48 GB 微调 / 推理};
 \draw[flow] (a1)--(a2);\draw[flow] (a2)--(a3);\draw[flow] (a3)--(a4);\draw[flow] (a4)--(a5);
 \node[lt] (l2) at (0,-2.75) {\faHospital\enspace 合作阶段 · 医院与生物合作者};
 \node[st,anchor=west] (b1) at (0,-3.7) {\ico{\faFileSignature}\ \textbf{审批与授权}\\伦理批件、人遗、协议};
 \node[st,right=.3cm of b1] (b2) {\ico{\faVials}\ \textbf{样本与测序}\\扰动实验需湿实验};
 \node[st,right=.3cm of b2] (b3) {\ico{\faServer}\ \textbf{算力位置}\\院内 GPU 或脱敏导出};
 \node[st,right=.3cm of b3,fill=PaleTeal,draw=DarkBlue,line width=.9pt] (b4) {\ico{\faLock}\ \textbf{同一套评分脚本}\\医院数据只用于评估};
 \node[st,right=.3cm of b4] (b5) {\ico{\faUserMd}\ \textbf{临床 / 生物核验}\\差异基因与通路解释};
 \draw[flow] (b1)--(b2);\draw[flow] (b2)--(b3);\draw[flow] (b3)--(b4);\draw[flow] (b4)--(b5);
 \draw[flow,dashed,{Latex[length=2mm]}-{Latex[length=2mm]}] (a4.south) -- node[right,font=\scriptsize,text=Muted]{复用} (b4.north);
 \begin{scope}[on background layer]
  \node[lane,fit=(l1)(a1)(a5)] {};
  \node[lane,fill=PaleAmber,fit=(l2)(b1)(b5)] {};
 \end{scope}
\end{tikzpicture}
\remark{\ico{\faHourglassHalf}\enspace 合作阶段的审批与准备以月计，应与起步阶段并行推进；起步阶段的划分、基线和评分脚本直接复用。}
\source{显存与数据规模为量级估计，需按所选模型与院方条件核对；合作前的关卡见下一页。}
\end{frame}
